
Soit $\beta_y:A\to K_\bullet(X)$ l'homomorphisme ainsi défini. Si l'on désigne par $I$ l'idéal $\Filt^1(X)/\Filt^{j+1}(X)$ de $A$, on a $I^{j+1}=0$. Gardant les notations du lemme 2.1.4 on trouve un homomorphisme de groupes abéliens
\[
\gamma:\Pic(X)\longrightarrow P
\]
qui associe à la classe de tout $\underline{O}_X$-module inversible $\underline{L}$ dans $\Pic(X)$ l'élément $j(\cl^{\bullet}(\underline{L}))$ de $P$.

On tire immédiatement des lemmes 2.1.2 et 2.1.4 :

\textbf{Proposition 2.2.1.} \uline{Pour tout schéma noethérien} $X$ \uline{et pour tout} $y\in\Filt_j(X)$ ($j\in\mathbb Z$\uline{) l'application composée}
\[
\Pic(X)\longrightarrow K_\bullet(X)\otimes_{\mathbb Z}\mathbb Q,
\]
\uline{définie par}
\[
\cl_P(\underline{L})\rightsquigarrow (\cl^{\bullet}(\underline{L})y)\otimes1
\]
\uline{est une application polynôme de degré} $\leq j$.

\subsection*{2.3.}
Si $X$ est un schéma propre sur un corps $k$ et si $g:X\to\Spec(k)$ est le morphisme structural, alors l'homomorphisme
\[
g_*:K_\bullet(X)\longrightarrow K_\bullet(\Spec(k))\cong\mathbb Z
\]
associe à la classe $\cl_\bullet(\underline{F})$ de tout $\underline{O}_X$-module cohérent $\underline{F}$ la caractéristique d'Euler-Poincaré $\chi(\underline{F})\in\mathbb Z$ de $\underline{F}$ sur $X$. De 2.2.1 on conclut immédiatement :

